\chapter{Barisan}\label{ch:g12:seq}

Sebuah \omterm{def:g12:seq:sequence}{barisan} adalah daftar \omterm{def:g10:numbers:sets}{bilangan real} yang terindeks oleh \omterm{def:g10:numbers:sets}{bilangan cacah}.
\omterm{def:g12:seq:sequence}{Barisan} memodelkan perkembangan diskret --- populasi yang dibilang tahun demi
tahun, saldo sebuah rekening bank, \omterm{def:g10:numbers:approx}{hampiran} sebuah bilangan yang berturut-turut
--- dan limitnya adalah perjumpaan serius pertama dengan ketakhinggaan. Bab ini
menyiapkan kosakatanya, asas induksi matematika, serta teorema kekonvergenan
yang mendasar.

\section{Bernalar dengan induksi matematika}

\begin{theorem}[Asas induksi matematika]\label{thm:g12:seq:induction}
Misalkan $P(n)$ sebuah pernyataan yang bergantung pada \omterm{def:g10:numbers:sets}{bilangan bulat} $n$, dan
misalkan $n_0 \in \N$. Jika
\begin{enumerate}
  \item \emph{(langkah dasar)} $P(n_0)$ benar, dan
  \item \emph{(langkah induksi)} untuk setiap $n \geq n_0$, $P(n)$ mengakibatkan $P(n+1)$,
\end{enumerate}
maka $P(n)$ benar untuk setiap $n \geq n_0$.
\end{theorem}

\begin{proof}
Andaikan, untuk memperoleh kontradiksi, bahwa himpunan $A$ berisi \omterm{def:g10:numbers:sets}{bilangan bulat}
$n \geq n_0$ yang membuat $P(n)$ salah tidaklah kosong. Maka $A$ mempunyai
anggota terkecil $m$.\footnote{Setiap himpunan bagian taktakosong dari $\N$
mempunyai anggota terkecil; sifat $\N$ ini diambil sebagai aksioma.} Karena
$P(n_0)$ benar, $m > n_0$, sehingga $m - 1 \geq n_0$ dan $m-1 \notin A$, yaitu
$P(m-1)$ benar. Langkah induksi yang diterapkan pada $n = m-1$ lalu
memperlihatkan bahwa $P(m)$ benar, yang bertentangan dengan $m \in A$.
\end{proof}

\begin{example}\label{ex:g12:seq:bernoulli}
Mari kita buktikan \emph{ketaksamaan Bernoulli}: untuk setiap \omterm{def:g10:numbers:sets}{bilangan real}
$a > 0$ dan setiap $n \in \N$,
\[
(1+a)^n \geq 1 + na.
\]
\emph{Langkah dasar.} Untuk $n = 0$, kedua ruasnya sama dengan $1$.
\emph{Langkah induksi.} Andaikan $(1+a)^n \geq 1+na$ untuk suatu $n \in \N$.
Karena $1 + a > 0$, mengalikan kedua ruasnya dengan $1+a$ mempertahankan ketaksamaannya:
\[
(1+a)^{n+1} \geq (1+na)(1+a) = 1 + (n+1)a + na^2 \geq 1 + (n+1)a .
\]
Menurut induksi, ketaksamaannya berlaku untuk semua $n \in \N$.
\end{example}

\begin{method}[Menulis bukti dengan induksi]\label{met:g12:seq:induction}
Nyatakan selalu pernyataan $P(n)$ secara tersurat sebelum mulai. Bukti yang
lengkap mempunyai tiga bagian yang terlihat: langkah dasarnya, langkah
induksinya (``andaikan $P(n)$; kita buktikan $P(n+1)$''), dan kesimpulan yang
memanggil asas induksi matematika. Kekeliruan yang paling lazim adalah
membuktikan langkah induksinya tanpa pernah memakai hipotesis $P(n)$: bila itu
terjadi, entah buktinya keliru entah induksinya memang tidak diperlukan.
\end{method}

\section{Kosakata barisan}

\begin{definition}[Barisan]\label{def:g12:seq:sequence}
Sebuah \emph{barisan}\index{barisan} adalah \omterm{def:g11:func:function}{fungsi} $u \colon \N \to \R$ (atau
dari $\{n \in \N : n \geq n_0\}$ ke $\R$). \omterm{def:g10:functions:function}{Peta} $n$ ditulis $u_n$, dan
barisannya sendiri ditulis $(u_n)_{n\in\N}$ atau cukup $(u_n)$.
\end{definition}

Sebuah \omterm{def:g12:seq:sequence}{barisan} dapat diberikan secara \emph{eksplisit}, lewat rumus
$u_n = f(n)$, atau lewat \emph{rekurensi}, lewat suku pertamanya dan hubungan
$u_{n+1} = f(u_n)$.

\begin{definition}[Kemonotonan]\label{def:g12:seq:monotonic}
Sebuah \omterm{def:g12:seq:sequence}{barisan} $(u_n)$ disebut \emph{naik}\index{barisan!naik} bila
$u_{n+1} \geq u_n$ untuk semua $n$, \emph{turun} bila $u_{n+1} \leq u_n$ untuk
semua $n$, dan \emph{monoton} bila \omterm{def:g12:seq:sequence}{barisan} itu naik atau turun. \omterm{def:g12:seq:sequence}{Barisan} itu
disebut \emph{tegas} naik (masing-masing turun) apabila ketaksamaannya
tegas.
\end{definition}

\begin{method}[Menelaah kemonotonan sebuah barisan]\label{met:g12:seq:monotonicity}
Tiga teknik baku:
\begin{enumerate}
  \item telaahlah tanda $u_{n+1} - u_n$;
  \item bila semua sukunya positif, bandingkan $\dfrac{u_{n+1}}{u_n}$ dengan $1$;
  \item bila $u_n = f(n)$ dengan $f$ terdefinisi pada $\intco{0}{+\infty}$,
        pakailah variasi $f$.
\end{enumerate}
\end{method}

\begin{definition}[Barisan terbatas]\label{def:g12:seq:bounded}
Sebuah \omterm{def:g12:seq:sequence}{barisan} $(u_n)$ disebut \emph{terbatas di atas}\index{barisan!terbatas}
bila ada $M \in \R$ dengan $u_n \leq M$ untuk semua $n$; \emph{terbatas di
bawah} bila ada $m \in \R$ dengan $u_n \geq m$ untuk semua $n$; dan
\emph{terbatas} bila keduanya berlaku.
\end{definition}

\subsection{Barisan aritmetika dan geometri}

\begin{definition}[Barisan aritmetika dan geometri]\label{def:g12:seq:arith-geom}
Sebuah \omterm{def:g12:seq:sequence}{barisan} $(u_n)$ disebut \emph{aritmetika}\index{barisan!aritmetika}
dengan \omterm{def:g11:seq:arithmetic}{beda} $r$ bila $u_{n+1} = u_n + r$ untuk semua $n$, dan
\emph{geometri}\index{barisan!geometri} dengan \omterm{def:g11:seq:geometric}{rasio} $q$ bila
$u_{n+1} = q\,u_n$ untuk semua $n$.
\end{definition}

\begin{proposition}[Bentuk eksplisit dan jumlahnya]\label{prop:g12:seq:arith-geom}
Misalkan $n \in \N$.
\begin{enumerate}
  \item Jika $(u_n)$ \omterm{def:g12:seq:arith-geom}{aritmetika} dengan \omterm{def:g11:seq:arithmetic}{beda} $r$, maka
        $u_n = u_0 + nr$ dan
        \[
        u_0 + u_1 + \dots + u_n = (n+1)\,\frac{u_0 + u_n}{2}.
        \]
  \item Jika $(u_n)$ \omterm{def:g12:seq:arith-geom}{geometri} dengan \omterm{def:g11:seq:geometric}{rasio} $q \neq 1$, maka
        $u_n = u_0\, q^n$ dan
        \[
        u_0 + u_1 + \dots + u_n = u_0\,\frac{1 - q^{n+1}}{1 - q}.
        \]
\end{enumerate}
\end{proposition}

\begin{proof}
Bentuk eksplisitnya menyusul lewat induksi langsung. Untuk jumlah
\omterm{def:g12:seq:arith-geom}{aritmetikanya}, tulislah $S = u_0 + \dots + u_n$ lalu tambahkan jumlah yang sama
dalam urutan terbalik: masing-masing dari $n+1$ jumlah kolomnya sama dengan
$u_0 + u_n$, sehingga $2S = (n+1)(u_0+u_n)$. Untuk jumlah \omterm{def:g12:seq:arith-geom}{geometrinya}, hitunglah
$S - qS$: semua sukunya saling menghapus berpasangan kecuali yang pertama dan
yang terakhir, sehingga $(1-q)S = u_0(1 - q^{n+1})$.
\end{proof}

\section{Limit sebuah barisan}

\begin{definition}[Barisan konvergen]\label{def:g12:seq:limit}
Sebuah \omterm{def:g12:seq:sequence}{barisan} $(u_n)$ \emph{konvergen}\index{limit!barisan} ke \omterm{def:g10:numbers:sets}{bilangan real}
$\ell$ bila setiap \omterm{def:g10:numbers:interval}{interval} terbuka yang memuat $\ell$ memuat semua suku
$u_n$ mulai dari suatu indeks. Kita lalu tulis $\lim\limits_{n\to+\infty} u_n = \ell$.

Setara dengan itu: untuk setiap $\varepsilon > 0$, ada $N \in \N$ sehingga
untuk semua $n \geq N$ berlaku $\abs{u_n - \ell} \leq \varepsilon$.
\end{definition}

\begin{omfigure}
\begin{tikzpicture}
\begin{axis}[omaxis,
  width=11cm, height=5.6cm,
  xlabel={$n$}, ylabel={$u_n$},
  xmin=0, xmax=16.5, ymin=0.7, ymax=2.9,
  xtick={4}, xticklabels={$N$}, ytick={2}, yticklabels={$\ell$}]
  \fill[omDef!20] (axis cs:0,1.75) rectangle (axis cs:16.5,2.25);
  \addplot[omProp, dashed, domain=0:16.5] {2};
  \addplot[omDef, only marks, mark=*, mark size=1.5pt,
    samples at={1,...,16}] {2 + (-1)^x/x};
  \draw[black, dashed, thin] (axis cs:4,0.7) -- (axis cs:4,2.25);
  \node[font=\small, anchor=east] at (axis cs:16.3,2.4) {$\ell+\varepsilon$};
  \node[font=\small, anchor=east] at (axis cs:16.3,1.6) {$\ell-\varepsilon$};
\end{axis}
\end{tikzpicture}

{\small Kekonvergenan $u_n = 2 + \frac{(-1)^n}{n}$ ke $\ell = 2$: diberikan
$\varepsilon > 0$, semua suku mulai dari indeks $N$ terletak dalam pita
$\intcc{\ell-\varepsilon}{\ell+\varepsilon}$.}
\end{omfigure}

\begin{definition}[Divergen ke takhingga]\label{def:g12:seq:infinity}
\omterm{def:g12:seq:sequence}{Barisan} $(u_n)$ \emph{menuju $+\infty$} bila untuk setiap $A \in \R$ ada
$N \in \N$ sehingga $u_n \geq A$ untuk semua $n \geq N$. Kita tulis
$\lim\limits_{n\to+\infty} u_n = +\infty$; definisi
$\lim u_n = -\infty$ serupa. \omterm{def:g12:seq:sequence}{Barisan} yang tidak \omterm{def:g12:seq:limit}{konvergen} dikatakan
\emph{divergen}.
\end{definition}

\begin{remark}
Sebuah \omterm{def:g12:seq:sequence}{barisan} dapat divergen tanpa menuju $\pm\infty$: \omterm{def:g12:seq:sequence}{barisan}
$u_n = (-1)^n$ hanya mengambil nilai $1$ dan $-1$ dan tidak berlimit.
\end{remark}

\begin{proposition}[Ketunggalan limitnya]\label{prop:g12:seq:uniqueness}
Jika $(u_n)$ \omterm{def:g12:seq:limit}{konvergen}, limitnya tunggal.
\end{proposition}

\begin{proof}
Andaikan $u_n \to \ell$ dan $u_n \to \ell'$ dengan $\ell \neq \ell'$, katakan
$\ell < \ell'$. Ambil $\varepsilon = \frac{\ell' - \ell}{3} > 0$. Mulai dari
suatu indeks berlaku $\abs{u_n - \ell} \leq \varepsilon$ dan
$\abs{u_n - \ell'} \leq \varepsilon$, sehingga
\[
\ell' - \ell \leq \abs{\ell' - u_n} + \abs{u_n - \ell}
\leq 2\varepsilon = \tfrac{2}{3}(\ell' - \ell) < \ell' - \ell,
\]
sebuah kontradiksi.
\end{proof}

\begin{proposition}[Operasi pada limit]\label{prop:g12:seq:operations}
Misalkan $(u_n)$ dan $(v_n)$ dua \omterm{def:g12:seq:sequence}{barisan} dengan limit $\ell$ dan $\ell'$
(berhingga atau takhingga). Maka, selama ruas kanannya bukan bentuk taktentu,
\[
\lim (u_n + v_n) = \ell + \ell', \qquad
\lim (u_n v_n) = \ell\,\ell', \qquad
\lim \frac{u_n}{v_n} = \frac{\ell}{\ell'}.
\]
Bentuk taktentunya adalah
$(+\infty) + (-\infty)$, $0 \times \infty$, $\frac{\infty}{\infty}$ dan
$\frac{0}{0}$.
\end{proposition}

\begin{proof}
Kita buktikan aturan jumlahnya untuk limit yang berhingga; kasus lainnya serupa
dan ditinggalkan sebagai latihan. Misalkan $\varepsilon > 0$. Ada $N_1, N_2$
sehingga $\abs{u_n - \ell} \leq \varepsilon/2$ untuk $n \geq N_1$ dan
$\abs{v_n - \ell'} \leq \varepsilon/2$ untuk $n \geq N_2$. Untuk
$n \geq \max(N_1, N_2)$, ketaksamaan segitiga memberi
\[
\abs{(u_n + v_n) - (\ell + \ell')}
\leq \abs{u_n - \ell} + \abs{v_n - \ell'} \leq \varepsilon. \qedhere
\]
\end{proof}

\begin{method}[Mengangkat bentuk taktentu]\label{met:g12:seq:indeterminate}
Ketika berhadapan dengan bentuk taktentu, keluarkan suku yang menguasainya sebagai faktor. Misalnya
\[
n^2 - n = n^2\left(1 - \tfrac{1}{n}\right) \xrightarrow[n\to+\infty]{} +\infty,
\qquad
\frac{2n^2+1}{n^2 - n} = \frac{2 + 1/n^2}{1 - 1/n}
\xrightarrow[n\to+\infty]{} 2 .
\]
\end{method}

\section{Teorema kekonvergenan}

\begin{theorem}[Teorema pembandingan dan teorema apit]\label{thm:g12:seq:squeeze}
Misalkan $(u_n)$, $(v_n)$, $(w_n)$ \omterm{def:g12:seq:sequence}{barisan}.
\begin{enumerate}
  \item Jika $u_n \leq v_n$ mulai dari suatu indeks dan $u_n \to +\infty$, maka
        $v_n \to +\infty$.
  \item \emph{(Teorema apit)}\index{teorema apit} Jika
        $u_n \leq v_n \leq w_n$ mulai dari suatu indeks dan $(u_n)$ serta $(w_n)$
        sama-sama \omterm{def:g12:seq:limit}{konvergen} ke limit yang sama $\ell$, maka $(v_n)$ \omterm{def:g12:seq:limit}{konvergen}
        ke~$\ell$.
\end{enumerate}
\end{theorem}

\begin{proof}
\emph{1.} Misalkan $A \in \R$. Karena $u_n \to +\infty$, ada $N$ dengan
$u_n \geq A$ untuk $n \geq N$; dengan memperbesar $N$ bila perlu,
$v_n \geq u_n \geq A$ untuk $n \geq N$.

\emph{2.} Misalkan $\varepsilon > 0$. Mulai dari suatu indeks, berlaku
$\ell - \varepsilon \leq u_n$ dan $w_n \leq \ell + \varepsilon$, sehingga
$\ell - \varepsilon \leq u_n \leq v_n \leq w_n \leq \ell + \varepsilon$, yaitu
$\abs{v_n - \ell} \leq \varepsilon$.
\end{proof}

\begin{example}
Untuk semua $n \geq 1$, $-\frac{1}{n} \leq \frac{(-1)^n}{n} \leq \frac{1}{n}$,
dan kedua batasnya menuju $0$; jadi $\frac{(-1)^n}{n} \to 0$.
\end{example}

\begin{theorem}[Teorema kekonvergenan monoton]\label{thm:g12:seq:monotone}
\omterm{def:g12:seq:sequence}{Barisan} \omterm{def:g12:seq:monotonic}{naik} yang \omterm{def:g12:seq:bounded}{terbatas di atas} bersifat \omterm{def:g12:seq:limit}{konvergen}. \omterm{def:g12:seq:sequence}{Barisan} turun
yang \omterm{def:g12:seq:bounded}{terbatas di bawah} bersifat \omterm{def:g12:seq:limit}{konvergen}. \omterm{def:g12:seq:sequence}{Barisan} \omterm{def:g12:seq:monotonic}{naik} yang tidak \omterm{def:g12:seq:bounded}{terbatas
di atas} menuju $+\infty$.
\end{theorem}

\begin{proof}[Bukti sebagian]
Kita buktikan pernyataan ketiganya. Misalkan $(u_n)$ \omterm{def:g12:seq:monotonic}{naik} dan tidak \omterm{def:g12:seq:bounded}{terbatas
di atas}, dan misalkan $A \in \R$. Karena $A$ bukan batas atas, ada $N$
dengan $u_N \geq A$; menurut kemonotonannya, $u_n \geq u_N \geq A$ untuk semua
$n \geq N$. Jadi $u_n \to +\infty$.

Kedua pernyataan kekonvergenannya bersandar pada sifat batas atas terkecil
$\R$; keduanya \emph{diterima tanpa bukti pada tingkat ini} (dan dibuktikan
pada tahun pertama universitas).
\end{proof}

\begin{remark}
Teoremanya menjamin \emph{keberadaan} limitnya tetapi tidak memberikan
nilainya. \omterm{def:g12:seq:sequence}{Barisan} \omterm{def:g12:seq:monotonic}{naik} yang \omterm{def:g12:seq:bounded}{terbatas di atas} oleh $M$ \omterm{def:g12:seq:limit}{konvergen} ke suatu
$\ell \leq M$, tidak harus ke $M$.
\end{remark}

\begin{theorem}[Limit barisan geometri]\label{thm:g12:seq:geometric}
Misalkan $q \in \R$.
\begin{enumerate}
  \item Jika $q > 1$, maka $q^n \to +\infty$.
  \item Jika $q = 1$, maka $q^n \to 1$.
  \item Jika $\abs{q} < 1$, maka $q^n \to 0$.
  \item Jika $q \leq -1$, maka $(q^n)$ divergen dan tidak berlimit.
\end{enumerate}
\end{theorem}

\begin{proof}
\emph{1.} Tulislah $q = 1 + a$ dengan $a > 0$. Ketaksamaan Bernoulli
(\cref{ex:g12:seq:bernoulli}) memberi $q^n \geq 1 + na \to +\infty$, lalu kita
simpulkan lewat pembandingan (\cref{thm:g12:seq:squeeze}).

\emph{2.} Langsung.

\emph{3.} Jika $q = 0$ pernyataannya jelas. Bila tidak, $\abs{q} < 1$ memberi
$1/\abs{q} > 1$, sehingga $(1/\abs{q})^n \to +\infty$ menurut butir 1, jadi
$\abs{q}^n \to 0$, dan $-\abs{q}^n \leq q^n \leq \abs{q}^n$ memungkinkan kita
menyimpulkannya lewat teorema apit.

\emph{4.} Untuk $q \leq -1$, $(q^{2n})$ mengambil nilai $\geq 1$ sedangkan
$(q^{2n+1})$ mengambil nilai $\leq -1$: tidak ada satu limit pun yang dapat
menarik kedua subbarisan itu.
\end{proof}

\begin{omfigure}
\begin{tikzpicture}
\begin{axis}[omaxis,
  width=10cm, height=5.6cm,
  xlabel={$n$}, ylabel={$q^n$},
  xmin=-0.5, xmax=13, ymin=-1.3, ymax=4.3,
  xtick={5,10}, ytick={-1,1}]
  \addplot[omThm, only marks, mark=*, mark size=1.5pt,
    samples at={0,...,6}] {1.25^x};
  \addplot[omDef, only marks, mark=*, mark size=1.5pt,
    samples at={0,...,12}] {0.75^x};
  \addplot[omProp, only marks, mark=*, mark size=1.5pt,
    samples at={0,...,12}] {(-0.8)^x};
  \node[omThm, font=\small, anchor=west] at (axis cs:6.3,3.7) {$q=1.25$};
  \node[omDef, font=\small, anchor=south west] at (axis cs:3.4,0.55) {$q=0.75$};
  \node[omProp, font=\small, anchor=north] at (axis cs:7,-0.75) {$q=-0.8$};
\end{axis}
\end{tikzpicture}

{\small Ketiga perilaku $(q^n)$: divergen ke $+\infty$ untuk
$q > 1$ (merah), \omterm{def:g12:seq:limit}{konvergen} ke $0$ untuk $\abs q < 1$ (biru), dan osilasi
teredam --- yang tetap \omterm{def:g12:seq:limit}{konvergen} ke $0$ --- untuk $-1 < q < 0$ (jingga).}
\end{omfigure}

\begin{method}[Barisan rekuren $u_{n+1} = f(u_n)$]\label{met:g12:seq:recurrent}
Untuk menelaah \omterm{def:g12:seq:sequence}{barisan} yang ditentukan oleh $u_{n+1} = f(u_n)$:
\begin{enumerate}
  \item buktikan dengan induksi bahwa $(u_n)$ tetap berada pada sebuah \omterm{def:g10:numbers:interval}{interval}
        $I$ tempat $f$ berperilaku baik (dan, sering kali, bahwa $(u_n)$ \omterm{def:g12:seq:monotonic}{monoton});
  \item simpulkan kekonvergenannya dari teorema kekonvergenan \omterm{def:g12:seq:monotonic}{monoton};
  \item terapkan limitnya pada hubungan $u_{n+1} = f(u_n)$: jika $f$
        kontinu dan $u_n \to \ell \in I$, maka $\ell$ memenuhi
        $f(\ell) = \ell$ (lihat \cref{ch:g12:limcont}); selesaikan \omterm{def:g10:algebra:equation}{persamaan} ini
        lalu pilihlah \omterm{def:g11:quad:discriminant}{akar} yang tepat.
\end{enumerate}
\end{method}

\begin{omfigure}
\begin{tikzpicture}
\begin{axis}[omaxis,
  width=8.4cm, height=6.4cm,
  xmin=-0.4, xmax=3.1, ymin=-0.4, ymax=2.7,
  xtick={0,1.414,1.848}, xticklabels={$u_0$,$u_1$,$u_2$}, ytick=\empty]
  \addplot[omProp, dashed, domain=-0.3:2.6] {x};
  \addplot[omDef, thick, domain=-0.4:2.9] {sqrt(x+2)};
  \draw[omThm, thick] (axis cs:0,0) -- (axis cs:0,1.414)
    -- (axis cs:1.414,1.414) -- (axis cs:1.414,1.848)
    -- (axis cs:1.848,1.848) -- (axis cs:1.848,1.961)
    -- (axis cs:1.961,1.961);
  \draw[black, dashed, thin] (axis cs:1.414,0) -- (axis cs:1.414,1.414);
  \draw[black, dashed, thin] (axis cs:1.848,0) -- (axis cs:1.848,1.848);
  \fill (axis cs:2,2) circle (1.5pt)
    node[below right, font=\small] {$\ell = 2$};
  \node[omDef, font=\small, above left] at (axis cs:2.75,2.18)
    {$y = f(x)$};
  \node[omProp, font=\small, anchor=north west] at (axis cs:1.08,0.82)
    {$y = x$};
\end{axis}
\end{tikzpicture}

{\small Lukisan tangga untuk $u_{n+1} = \sqrt{u_n + 2}$, $u_0 = 0$
(\cref{exo:g12:seq:6}): setiap langkah tegaknya membaca $f(u_n)$ pada kurvanya,
setiap langkah mendatarnya membawanya kembali lewat $y = x$. \omterm{def:g12:seq:sequence}{Barisannya} mendaki
ke \omterm{pb:g10:functions:1}{titik tetap} $\ell = 2$, tempat kurvanya bertemu garisnya.}
\end{omfigure}

\begin{example}\label{ex:g12:seq:sqrt2}
Misalkan $u_0 = 2$ dan $u_{n+1} = \frac{1}{2}\left(u_n + \frac{2}{u_n}\right)$.
Dengan induksi diperiksa bahwa $u_n \geq \sqrt{2}$ untuk semua $n$ (ketaksamaan
$\frac{1}{2}(x + 2/x) \geq \sqrt{2}$ untuk $x>0$ setara dengan
$(x - \sqrt2)^2 \geq 0$), lalu bahwa $(u_n)$ turun, sebab
\[
u_{n+1}-u_n=\frac{2-u_n^2}{2u_n}\leq 0 .
\]
Karena turun dan \omterm{def:g12:seq:bounded}{terbatas di bawah}, $(u_n)$
\omterm{def:g12:seq:limit}{konvergen} ke suatu $\ell \geq \sqrt{2}$, yang pasti memenuhi
$\ell = \frac{1}{2}(\ell + 2/\ell)$, yaitu $\ell^2 = 2$. Jadi
$u_n \to \sqrt{2}$. Inilah algoritme Heron, yang sudah dipakai orang
Babilonia; kekonvergenannya sangat cepat ($u_3$ sudah memberi $\sqrt 2$
sampai delapan angka desimal).
\end{example}

\section{Latihan}

\begin{exercise}[$\star$]\label{exo:g12:seq:1}
Buktikan dengan induksi bahwa untuk semua $n \in \N$,
\[
1^2 + 2^2 + \dots + n^2 = \frac{n(n+1)(2n+1)}{6}.
\]
\end{exercise}

\begin{exercise}[$\star$]\label{exo:g12:seq:2}
Telaahlah kemonotonan \omterm{def:g12:seq:sequence}{barisan} yang ditentukan untuk $n \geq 1$ oleh
\[
a_n = \frac{n+1}{n}, \qquad
b_n = \frac{2^n}{n}, \qquad
c_n = n^2 - 10n .
\]
\end{exercise}

\begin{exercise}[$\star$]\label{exo:g12:seq:3}
Hitunglah limit \omterm{def:g12:seq:sequence}{barisan} yang suku umumnya
\[
u_n = \frac{3n^2 - n + 1}{2n^2 + 5}, \qquad
v_n = \sqrt{n+1} - \sqrt{n}, \qquad
w_n = \frac{2^n - 3^n}{3^n + 1}.
\]
\end{exercise}

\begin{exercise}[$\star$]\label{exo:g12:seq:4}
Misalkan $(u_n)$ \omterm{def:g11:seq:arithmetic}{barisan aritmetika} dengan $u_0 = 5$ dan \omterm{def:g11:seq:arithmetic}{beda}
$r = 3$, dan $(v_n)$ \omterm{def:g11:seq:geometric}{barisan geometri} dengan $v_0 = 8$ dan \omterm{def:g11:seq:geometric}{rasio}
$q = \frac{1}{2}$. Hitunglah $u_n$, $v_n$, $\sum_{k=0}^{n} u_k$ dan
$\sum_{k=0}^{n} v_k$, serta limit keempat bentuk itu ketika
$n \to +\infty$.
\end{exercise}

\begin{exercise}[$\star\star$]\label{exo:g12:seq:5}
Dengan teorema apit, hitunglah
\[
\lim_{n\to+\infty} \frac{n + \cos n}{n + 1}
\qquad\text{dan}\qquad
\lim_{n\to+\infty} \frac{n!}{n^n},
\]
dengan $n! = 1 \times 2 \times \dots \times n$. Untuk limit yang kedua, batasi
$\frac{n!}{n^n}$ oleh suku sebuah \omterm{def:g11:seq:geometric}{barisan geometri}.
\end{exercise}

\begin{exercise}[$\star\star$]\label{exo:g12:seq:6}
Misalkan $u_0 = 0$ dan $u_{n+1} = \sqrt{u_n + 2}$ untuk semua $n \in \N$.
\begin{enumerate}
  \item Buktikan dengan induksi bahwa $0 \leq u_n \leq 2$ untuk semua $n$.
  \item Tunjukkan bahwa $(u_n)$ \omterm{def:g12:seq:monotonic}{naik}.
  \item Simpulkan bahwa $(u_n)$ \omterm{def:g12:seq:limit}{konvergen} lalu tentukan limitnya.
\end{enumerate}
\end{exercise}

\begin{exercise}[$\star\star$]\label{exo:g12:seq:7}
Seorang pasien meminum dosis $1$ satuan obat setiap pagi. Selama setiap
selang 24 jam, tubuhnya menyingkirkan $40\%$ obat yang ada. Misalkan $u_n$
banyaknya obat dalam tubuh tepat setelah dosis pada hari ke-$n$, sehingga
$u_0 = 1$.
\begin{enumerate}
  \item Berikan alasan bahwa $u_{n+1} = 0.6\,u_n + 1$.
  \item Misalkan $v_n = u_n - 2.5$. Tunjukkan bahwa $(v_n)$ \omterm{def:g12:seq:arith-geom}{geometri} lalu
        simpulkan rumus eksplisit untuk $u_n$.
  \item Tentukan banyaknya obat dalam tubuh dalam jangka panjang.
\end{enumerate}
\end{exercise}

\begin{exercise}[$\star\star$]\label{exo:g12:seq:8}
Misalkan $(u_n)$ ditentukan oleh $u_0 = 3$ dan $u_{n+1} = \frac{4u_n - 1}{u_n + 2}$.
\begin{enumerate}
  \item Tunjukkan dengan induksi bahwa $u_n > 1$ untuk semua $n \in \N$.
  \item Tunjukkan bahwa $v_n = \dfrac{1}{u_n - 1}$ menentukan \omterm{def:g11:seq:arithmetic}{barisan aritmetika}.
  \item Simpulkan rumus eksplisit untuk $v_n$ dan $u_n$, serta limit
        $(u_n)$.
\end{enumerate}
\end{exercise}

\begin{exercise}[$\star\star\star$]\label{exo:g12:seq:9}
Untuk $n \geq 1$, misalkan
$H_n = 1 + \frac{1}{2} + \frac{1}{3} + \dots + \frac{1}{n}$.
\begin{enumerate}
  \item Tunjukkan bahwa untuk semua $n \geq 1$, $H_{2n} - H_n \geq \frac{1}{2}$.
  \item Simpulkan bahwa $H_{2^k} \geq 1 + \frac{k}{2}$ untuk semua $k \in \N$,
        lalu simpulkan bahwa $H_n \to +\infty$.
\end{enumerate}
\end{exercise}

\begin{exercise}[$\star\star\star$]\label{exo:g12:seq:10}
\emph{(\omterm{def:g12:seq:sequence}{Barisan} berdampingan.)} Dua \omterm{def:g12:seq:sequence}{barisan} $(a_n)$ dan $(b_n)$ disebut
\emph{berdampingan} bila $(a_n)$ \omterm{def:g12:seq:monotonic}{naik}, $(b_n)$ turun, dan
$b_n - a_n \to 0$.
\begin{enumerate}
  \item Tunjukkan bahwa untuk semua $n$ berlaku $a_n \leq b_n$. (Petunjuk:
        telaahlah kemonotonan $(b_n - a_n)$.)
  \item Tunjukkan bahwa \omterm{def:g12:seq:sequence}{barisan} yang berdampingan sama-sama \omterm{def:g12:seq:limit}{konvergen}, ke limit
        yang sama.
  \item Penerapannya: tunjukkan bahwa \omterm{def:g12:seq:sequence}{barisan}
        $a_n = \sum_{k=0}^{n} \frac{1}{k!}$ dan
        $b_n = a_n + \frac{1}{n \cdot n!}$ ($n \geq 1$) berdampingan. (Limit
        bersamanya adalah bilangan $\eu$, yang ditelaah pada \cref{ch:g12:exp}.)
\end{enumerate}
\end{exercise}

\section{Soal: Barisan Heron, akhirnya diadili}

\begin{problem}[{Soal akhir pekan --- induksi mengesahkan, kekonvergenan
monoton menjatuhkan vonis, dan resep berumur dua ribu tahun untuk
$\sqrt2$ akhirnya memperoleh buktinya (dengan rataan menakjubkan milik
Gauss sebagai penutup)}]
\label{pb:g12:seq:1}
Tiga kali seri buku ini berjumpa dengan resep Heron --- rata-ratakan
terkaan dengan $2/\text{terkaan}$ --- dan tiga kali pula hanya dapat
\emph{mengamati} bahwa resep itu berhasil. Bab ini akhirnya memiliki
alat penghakimannya: induksi
(\cref{thm:g12:seq:induction}), teorema kekonvergenan \omterm{def:g12:seq:monotonic}{monoton}
(\cref{thm:g12:seq:monotone}), dan limit rekurensi. Vonisnya,
beserta laju yang tersahkan, menempati jantung soal
ini; di sekelilingnya, jebakan klasik induksi, kedivergenan paling lambat
dalam matematika, dan kekonvergenan tercepat yang pernah ditemukan
Gauss.

\medskip
\noindent\textbf{Bagian I --- Pemanasan induksi.}
\begin{enumerate}
  \item Buktikan dengan induksi:
        $1 + 3 + 5 + \dots + (2n - 1) = n^2$ (tangga
        bilangan ganjil, yang digambar di jilid sebelumnya, kini
        tersahkan).
  \item Buktikan dengan induksi bahwa $2^n > n$ untuk setiap
        $n \in \N$.
  \item Buktikan ketaksamaan Bernoulli dengan induksi: untuk
        $x \geq 0$ dan $n \in \N$, berlaku
        $(1 + x)^n \geq 1 + nx$.
  \item Jebakan klasiknya: ``semua kelereng berwarna sama ---
        benar untuk satu kelereng; dan jika sebarang $n$ kelereng selalu
        berwarna tunggal, maka di antara $n + 1$ kelereng, $n$ yang
        pertama sewarna, $n$ yang terakhir sewarna, sehingga seluruh
        $n + 1$ kelereng pun sewarna.'' Setiap anak tahu kesimpulan itu
        mustahil: carilah langkah persisnya tempat induksinya patah.
  \item Buktikan dengan induksi bahwa $4^n - 1$ habis dibagi $3$
        untuk setiap $n \in \N$.
\end{enumerate}

\medskip
\noindent\textbf{Bagian II --- Pengadilan Heron.}
Misalkan $x_0 = 2$ dan $x_{n+1} = \dfrac12\left(x_n +
\dfrac{2}{x_n}\right)$.
\begin{enumerate}[resume]
  \item Hitunglah $x_1$, $x_2$, $x_3$ sebagai pecahan eksak (kawan-kawan
        lama).
  \item Buktikan identitas kuncinya
        \[
        x_{n+1}^2 - 2 = \left(\frac{x_n^2 - 2}{2x_n}\right)^{\!2}
        \geq 0,
        \]
        lalu simpulkan dengan induksi bahwa $x_n > 0$ dan
        $x_n^2 > 2$ untuk setiap $n$.
  \item Tunjukkan bahwa $(x_n)$ tegas turun (hitunglah
        $x_{n+1} - x_n$ lalu pakai pertanyaan 7).
  \item Panggillah teorema kekonvergenan \omterm{def:g12:seq:monotonic}{monoton}: mengapa
        $(x_n)$ \omterm{def:g12:seq:limit}{konvergen} ke suatu limit $L \geq 1$?
  \item Kenali limitnya: terapkan limit pada rekurensinya
        (\cref{prop:g12:seq:operations}) lalu simpulkan
        $L = \sqrt2$. Nyatakan vonis sejarahnya: setelah dua
        ribu tahun mengabdi dengan setia, resep Heron
        \emph{terbukti} \omterm{def:g12:seq:limit}{konvergen}.
  \item Laju yang tersahkan: dengan $e_n = x_n - \sqrt2$, buktikan
        \[
        e_{n+1} = \frac{e_n^2}{2 x_n} ,
        \]
        lalu simpulkan $e_{n+1} \leq \frac{e_n^2}{2\sqrt2}$: jadi
        galatnya dikuadratkan pada setiap langkah --- penggandaan angka
        yang teramati sejak jilid sebelumnya, kini menjadi teorema.
  \item Benarkan secara numeris: hitunglah $e_0, e_1, e_2, e_3$ (dari
        pertanyaan 6) lalu periksalah bahwa setiap $\frac{e_{n+1}}{e_n^2}$
        dekat dengan $\frac{1}{2x_n}$.
\end{enumerate}

\medskip
\noindent\textbf{Bagian III --- Kedivergenan paling lambat.}
\begin{enumerate}[resume]
  \item \Cref{exo:g12:seq:9} membuktikan
        $H_{2^k} \geq 1 + \frac k2$ untuk jumlah harmoniknya. Berapa
        banyak suku yang menjamin $H_n > 10$? (Sebuah pangkat dua sudah
        cukup; kagumilah besarnya.)
  \item Sebagai bandingannya, jumlah \omterm{def:g12:seq:arith-geom}{geometri}
        $1 + \frac12 + \frac14 + \dots + \frac{1}{2^n} =
        2 - \frac{1}{2^n}$ \omterm{def:g12:seq:limit}{konvergen} ke $2$
        (\cref{thm:g12:seq:geometric}): gerak hati batang cokelat
        di jilid sebelumnya, akhirnya menjadi pernyataan limit.
        Tulislah bukti dua barisnya.
  \item Di antara keduanya: tunjukkan bahwa jumlah
        $S_n = 1 + \frac{1}{4} + \frac{1}{9} + \dots +
        \frac{1}{n^2}$ \omterm{def:g12:seq:limit}{konvergen}, dengan membatasi
        $\frac{1}{k^2} \leq \frac{1}{k(k-1)} =
        \frac{1}{k-1} - \frac{1}{k}$ (untuk $k \geq 2$),
        secara teleskopik, lalu menerapkan kekonvergenan \omterm{def:g12:seq:monotonic}{monoton}. (Limitnya,
        $\frac{\pi^2}{6}$, adalah salah satu keajaiban Euler,
        yang dibuktikan pada jilid universitas.)
  \item Nyatakan moral pertanyaan 13--15 dalam dua kalimat:
        apa yang diputuskan oleh ``sukunya menuju $0$'' tentang
        kekonvergenan jumlahnya --- dan apa yang tidak?
\end{enumerate}

\medskip
\noindent\textbf{Bagian IV --- Rataan aritmetika--geometri milik
Gauss.}
Misalkan $a_0 = 1$, $b_0 = 2$, dan
\[
a_{n+1} = \sqrt{a_n b_n},
\qquad
b_{n+1} = \frac{a_n + b_n}{2} .
\]
\begin{enumerate}[resume]
  \item Hitunglah $a_1, b_1, a_2, b_2$ (lima angka desimal). Apa yang
        kamu amati tentang lajunya?
  \item Tunjukkan bahwa $a_n \leq b_n$ untuk setiap $n$ (yaitu
        ketaksamaan rataan aritmetika--geometri, yang dijumpai di sepanjang
        seri ini), bahwa $(a_n)$ \omterm{def:g12:seq:monotonic}{naik} dan $(b_n)$ turun.
  \item Tunjukkan bahwa
        $b_{n+1} - a_{n+1} \leq \frac{b_n - a_n}{2}$
        (faktorkan $b_{n+1} - a_{n+1} =
        \frac{(\sqrt{b_n} - \sqrt{a_n})^2}{2}$ lalu bandingkan),
        lalu simpulkan dengan \cref{exo:g12:seq:10} bahwa kedua
        \omterm{def:g12:seq:sequence}{barisan} itu berdampingan: keduanya berbagi limit bersama
        $M(1, 2)$, yaitu \emph{rataan aritmetika--geometri}.
  \item Hitunglah $M(1, 2)$ sampai enam angka desimal (berapa iterasi
        yang kamu perlukan?). Pada 30 Mei 1799, Gauss menghitung
        $M(1, \sqrt2)$ sampai sebelas angka desimal, mengenali
        $\frac{\pi}{M(1,\sqrt2)}$ sebagai sebuah integral yang sudah
        dikenalnya, lalu menulis bahwa sebuah ``lapangan baru analisis''
        telah terbuka --- dan memang demikian: integral eliptik, yang
        diceritakan pada jilid universitas. Tutuplah dengan laju
        kekonvergenan yang teramati pada soal ini, dari yang paling lambat
        sampai yang paling cepat.

\end{enumerate}
\end{problem}
